problem:  
Let \( (M^4,g,f) \) be a compact smooth solution of the **Critical Point Equation (CPE)**  
\[
\mathrm{Ric} - \frac{R}{4} g = \nabla^2 f - \frac{\Delta f}{4} g,
\]
where \(R\) is the scalar curvature, \(f\) the associated potential, and \(M\) has **constant scalar curvature**.  
Assume further that \(g\) is **Bach-flat**, i.e. \( \mathrm{Bach}(g)=0\).  

It is known that Einstein metrics solve the CPE trivially (with \(f\) constant). Nontrivial CPE metrics remain mysterious, and no compact, non-Einstein, Bach-flat examples are known.

**Question:**  
Is it true that every compact, non-Einstein CPE metric with Bach-flat metric tensor must (up to finite covering) be *locally conformal* to a product of two constant curvature surfaces?  
Equivalently, does the combination of conditions  
\[
\text{CPE} + \text{Bach-flat} + R = \text{constant}
\]
force \( (M,g) \) to belong to the conformal class of a product geometry \( (\Sigma_1^2 \times \Sigma_2^2, g_1 + g_2) \)?  
If not, can one construct counterexamples, e.g. via warped-product deformations of Einstein metrics, or prove an obstruction result constraining the topology of such manifolds?

---

Why is it a "good" problem:  
This problem focuses on the intricate intersection of three fundamental structures in 4‑dimensional geometry:  

1. **Variational structure (CPE metrics):** These arise as critical points of the total scalar curvature functional under a volume constraint, linking curvature, potential theory, and variational analysis.  

2. **Conformal geometry (Bach-flatness):** The Bach tensor is the conformally invariant Euler–Lagrange tensor for the Weyl functional in dimension four—its vanishing means the metric sits at a critical point of a conformally invariant energy.  

3. **Global structure and classification:** Known classifications (notably results of Derdzinski, Chang–Gursky–Yang, etc.) describe harmonic-Weyl or half-flat situations, but the *CPE + Bach-flat* combination is unexplored. It demands connecting analysis of an elliptic PDE system with conformal geometry—to determine whether variational criticality enforces product‑type rigidity.  

The question stands at the frontier of current understanding because:  
- No non-Einstein compact CPE Bach-flat metrics are known.  
- Standard decomposition theorems do not directly apply to the CPE tensor.  
- A positive resolution or counterexample would reveal deep structural truths about the interaction between conformal and variational curvature conditions in 4‑dimensional Riemannian geometry.  

It is thus a *simple to state yet profoundly challenging* open classification problem, bridging geometric analysis, conformal geometry, and global differential topology.